\documentclass[11pt]{article}
\usepackage[utf8]{inputenc}
\usepackage{amsmath,amssymb}
\usepackage[margin=1in]{geometry}
\usepackage{hyperref}
\emergencystretch=3em

\title{The multiplicity theorem without ordering hypotheses:\\
$\int_{(0,b]^{D+2}}u^h f(\sqrt n\,u^k)\,du\sim C\,n^{-p/2}(\log n)^m$}
\author{Claude, with Daniel Murfet}
\date{2026-09-06}

\begin{document}
\maketitle
\sloppy

\begin{abstract}
For arbitrary exponent vectors $k,h:\mathrm{Fin}(D+2)\to\mathbb N$ with $k_i>0$, let $p$ be the
minimum of the candidate exponents $q_i=(h_i+1)/k_i$ and $m+1$ the number of $i$ attaining it. Then
the box integral at scale $\sqrt n$ is asymptotically $C\,n^{-p/2}(\log n)^m$ for some $C>0$
(\texttt{boxIntegralFin\_isEquivalent\_general}). The coordinates are sorted with
\texttt{Tuple.sort}, the first index attaining the minimum is a \texttt{Finset.min'}, the
multiplicity is computed through \texttt{Finset.card\_equiv} and \texttt{Fin.card\_Ici}, and the two
cases (all exponents equal; some strictly larger) are units~50 and~58. This closes the leading-term
theory of \texttt{thm:TaylorTree} for constant $\xi$ and monomial $\eta$ in the chart geometry.
Zero \texttt{sorry}/\texttt{axiom}.
\end{abstract}

\section{Setting}
Unit~58 needed a relabelling putting the exponents in the sorted shape. \texttt{Tuple.sort} applied
to $-q$ gives a permutation $\sigma$ with $q\circ\sigma$ antitone; then $p=q(\sigma(\mathrm{last}))$
is the minimum, the set of indices $j$ with $q(\sigma j)=p$ is an upper set, hence the interval
$[i_0,\mathrm{last}]$ for its least element $i_0$, and everything below $i_0$ is strictly above $p$.

\section{The things formalised}
{\small
\begin{verbatim}
noncomputable def finExp {d : ℕ} (k h : Fin d → ℕ) (i : Fin d) : ℝ := ((h i : ℝ) + 1) / k i
noncomputable def sortPerm {d : ℕ} (k h : Fin d → ℕ) : Equiv.Perm (Fin d) :=
  Tuple.sort fun i => -(finExp k h i)
theorem sortPerm_antitone {d : ℕ} (k h : Fin d → ℕ) : Antitone (finExp k h ∘ sortPerm k h)
theorem boxIntegralFin_isEquivalent_general (β a b : ℝ) (hβ : 0 < β) (hb : 0 < b) (D : ℕ)
    (k h : Fin (D + 2) → ℕ) (hk : ∀ i, 0 < k i) :
    ∃ (p : ℝ) (m : ℕ) (C : ℝ), 0 < C ∧ (∀ i, p ≤ finExp k h i) ∧
      m + 1 = (Finset.univ.filter fun i => finExp k h i = p).card ∧
      (fun n => boxIntegralFin β a b (Real.sqrt n) k h) ~[atTop]
        fun n => C * (n ^ (-(p / 2)) * Real.log n ^ m)
\end{verbatim}
}

\section{Proof strategy}
With $i_0$ the least index of the minimal set $S$: $S=\mathrm{Ici}\,i_0$ (antitone plus minimality),
so the multiplicity is $D+2-i_0$. If $i_0=0$ all exponents equal $p$ and unit~50 applies through the
bridge of units~54 and~58, with $m=D+1$ and constant $(\prod k_i^{-1})A_{p-1}/(2^{D+1}(D+1)!)$.
Otherwise $i_0=r+1$, $m=D+1-i_0$, and the sorted theorem in its $r+2+m=d$ form
(\texttt{boxIntegralFin\_isEquivalent\_of\_perm'}) applies with the sortedness, block and strictness
hypotheses read off through \texttt{chartExp\_toNatFun\_eq\_finExp}; the constant is positive because
each factor is (\texttt{genScale\_pos}, \texttt{genTower\_logBase}).

\section{Roadblocks and resolutions}
\begin{itemize}
\item \texttt{Tuple.sort} lives in \texttt{Mathlib.Data.Fin.Tuple.Sort}, not in the import closure.
\item \texttt{linarith} does not unfold the definition \texttt{sortPerm}; \texttt{simp only [sortPerm]}
first.
\item Membership of the last index in the minimal set is \texttt{rfl} after unfolding the filter;
\texttt{simp} stopped at the \texttt{set}-bound variable.
\item The unit~58 theorem fixed the dimension as \texttt{r + 2 + m} in the types; a variant with the
equation \texttt{r + 2 + m = d} (rewritten on the $\mathbb N$ side only) avoids transporting along a
propositional equality of types.
\end{itemize}

\section{What was Mathlib, what was new}
Mathlib: \texttt{Tuple.sort}, \texttt{Tuple.monotone\_sort}, \texttt{Finset.min'} and its API,
\texttt{Finset.card\_equiv}, \texttt{Fin.card\_Ici}, \texttt{Fin.le\_last}. New: the hypothesis-free
multiplicity theorem.

\section{Lessons}
Keep type-level dimensions fixed in the workhorse theorems and add a thin variant taking the dimension
as a propositional equation; sorting arguments then live entirely on the $\mathbb N$ side.

\section{Outcome}
Units~44--59 formalise the leading-term content of \texttt{thm:TaylorTree} in the chart geometry for
constant $\xi$ and monomial $\eta$: for every exponent vector the box integral is
$C\,n^{-p/2}(\log n)^m$ with $p$ half the minimal candidate exponent and $m+1$ its multiplicity, and
in the sorted form the constant and a complete polynomial-in-$\log n$ expansion are explicit.
\end{document}
